\section{端到端：从功能开发到能力训练}

上一章讲 BEV 如何统一空间，本章看更大的转向：端到端。郎咸朋说“我们以前做自动驾驶都做错了”，指的是把自动驾驶当作一个功能来开发：列出场景，写规则，划 ODD，逐个补丁修 bug。但真实驾驶场景几乎无法穷举，变量之间也不独立，改好一个场景可能破坏另一个场景。端到端的意义，是把驾驶从“人工写规则”转成“模型从数据中学习能力”。

\begin{figure}[H]
\centering
\includegraphics[width=0.96\textwidth]{e2e-stack.png}
\caption{端到端自动驾驶栈：从传感器输入到轨迹输出，中间层逐步被学习系统压缩。自制概念图，依据 00:28:46--00:41:14 对谈内容整理。}
\end{figure}

\begin{knowledgebox}{读图：端到端不是没有结构，而是让模型承担更多映射}
从 Sensors、BEV、Planning、Trajectory 到 Control，传统系统每层都有大量人工规则。端到端并不是完全消灭工程，而是把关键映射交给模型和数据，让系统学会从输入到轨迹的能力。
\end{knowledgebox}

\subsection{为什么场景穷举不可行}

本节拆解规则系统的失败。天气、车流、道路结构、路面状态、周围车辆、行人、自车速度、施工、坑洞、异物等变量组合巨大。即使人为定义上千万种场景，也会遇到两个问题：场景之间不正交，规则互相影响；长尾事件永远会出现，人类能用常识避开，规则系统却可能不知道如何处理。

\begin{warningbox}{规则系统的长尾困境}
自动驾驶不是按钮逻辑。一个右转规则在晴天、雨天、堵车、行人、自行车、施工和路面坑洞下会完全不同。规则越多，互相干扰越严重，系统维护成本越高。
\end{warningbox}

\subsection{能力系统：像教孩子，不是替孩子学}

本节保留郎咸朋非常关键的口头比喻。端到端之后，工程师像旁观者和老师：提供高质量数据、训练资源、评测题和反馈，但不能替模型学习。模型学成什么样，内部怎么形成策略，工程师未必完全可解释；但可以通过 case、回归测试、仿真、路测和用户反馈评价它是否进步。

\begin{importantbox}{课堂提示：模型和数据是同一枚硬币的两面}
用模型做能力系统，就必然依赖数据驱动；拥有高质量、垂直、真实反馈数据的公司，才有机会持续改进自动驾驶。模型参数、算法技巧和团队经验重要，但没有高质量驾驶数据，天花板会很低。
\end{importantbox}

\subsection{端到端之后，工程师还做什么}

本节补一个容易误解的点。端到端不是说工程师退场，也不是把安全责任交给黑盒。相反，工程师的工作从“写每个场景的规则”转向“定义数据、评测、回归和安全边界”。他们要找长尾 case，判断哪些场景暴露了模型盲区；要设计标注和训练数据，让模型看到高价值样本；要建立仿真和闭环评测，避免新模型修好一个场景又破坏另一个场景；还要控制上线节奏，确保量产用户不是未经保护的实验对象。

\begin{knowledgebox}{实践经验：端到端把工程难点搬到了数据和评测}
规则时代的工程难点在代码分支和场景穷举；模型时代的工程难点在数据覆盖、样本权重、训练稳定性、离线指标、在线体验和安全回归。不是工程量消失，而是工程量换了位置。
\end{knowledgebox}

\noindent\begin{tabularx}{\textwidth}{p{0.22\textwidth}p{0.34\textwidth}X}
\toprule
工程角色 & 规则时代主要工作 & 端到端时代主要工作 \\
\midrule
算法工程师 & 写检测、规划、规则和状态机 & 设计模型结构、训练目标、数据采样和评测。 \\
数据工程师 & 收集少量 case 验证规则 & 建立车队回流、清洗、脱敏、标注和挖掘管线。 \\
测试工程师 & 跑固定场景和手工回归 & 做长尾 case 库、仿真、离线指标和线上灰度。 \\
产品/安全 & 定义功能边界和提示 & 定义接管责任、ODD、风险降级和用户教育。 \\
\bottomrule
\end{tabularx}

\subsection{端到端、RL、VLM 与世界模型}

本节连接后续路线。访谈中把端到端、VLM、世界模型和 RL 放在一起讨论。端到端先让系统从感知到轨迹学习驾驶行为；VLM 引入语义理解和复杂场景解释；世界模型预测行动对未来交通状态的影响；RL 则提供闭环优化框架。自动驾驶越来越接近物理世界里的 Agent：观察、预测、行动、反馈、再学习。

\begin{figure}[H]
\centering
\includegraphics[width=0.96\textwidth]{e2e-rl-worldmodel.png}
\caption{端到端 + VLM + 世界模型：从模仿轨迹走向强化学习和未来状态预测。自制概念图，依据 01:15:15--01:26:40 对谈内容整理。}
\end{figure}

\begin{knowledgebox}{读图：世界模型补的是“后果预测”}
端到端策略可以输出轨迹，但复杂交通需要预测：如果我变道、减速、让行，周围车和人会怎样变化。世界模型把动作后果纳入系统，RL 则让系统在反馈中优化策略。
\end{knowledgebox}

\subsection{本章小结}

端到端的本质不是一个流行词，而是自动驾驶从软件功能开发转向能力训练。它把工程师从手写规则者变成数据、训练、评测和安全边界的设计者。

